%% ============================================================================
%%  back/glossary.tex
%%  واژه‌نامه‌های پایان‌نامه (توصیفی، فارسی به انگلیسی، انگلیسی به فارسی)
%%  طراحی شده برای قالب رسمی دانشگاه خوارزمی
%% ============================================================================
%%
%%  توضیح:
%%  این فایل شامل سه بخش تخصصی واژه‌نامه مطابق با جدول ۲-۱ شیوه‌نامه دانشگاه خوارزمی است:
%%  ۱. واژه‌نامه توصیفی: تعریف و تشریح دقیق اصطلاحات و مفاهیم کلیدی پژوهش
%%  ۲. واژه‌نامه فارسی به انگلیسی: معادل‌های زبان اصلی به ترتیب الفبای فارسی
%%  ۳. واژه‌نامه انگلیسی به فارسی: معادل‌های زبان فارسی به ترتیب الفبای انگلیسی
%% ============================================================================

% -----------------------------------------------------------------------------
% ۱. واژه‌نامه توصیفی
% -----------------------------------------------------------------------------
\begin{descriptiveglossary}
\item[یادگیری عمیق:] زیرشاخه‌ای نوین از یادگیری ماشین که مبتنی بر شبکه‌های عصبی چندلایه با ساختار سلسله‌مراتبی برای استخراج خودکار ویژگی‌های پیچیده از کلان‌داده‌ها عمل می‌کند.

\item[مدل ترنسفورمر:] معماری پردازش توالی مبتنی بر سازوکار توجه خودکار که ارتباطات دوربرد میان کلمات یا داده‌ها را بدون وابستگی به لایه‌های بازگشتی مدل‌سازی می‌نماید.

\item[سازوکار توجه:] تابعی ریاضی برای وزن‌دهی به بخش‌های گوناگون دنباله ورودی، به‌گونه‌ای که مدل بتواند بر مرتبط‌ترین اجزا در زمان تولید خروجی تمرکز کند.

\item[حروف‌چینی دیجیتال:] فرآیند چیدمان و صفحه‌آرایی دقیق متون، فرمول‌های ریاضی، جداول و عناصر گرافیکی با بهره‌گیری از موتورهای محاسباتی مانند تک و لاتک.

\item[زی‌پرشین:] بسته‌ای جامع در سامانه زی‌تک که امکان حروف‌چینی متن‌های فارسی و دوزبانه (راست‌به‌چپ و چپ‌به‌راست) را با حفظ کامل اصول تایپوگرافی فراهم می‌آورد.

\item[تابع زیان:] تابعی هدف در بهینه‌سازی که تفاوت میان مقادیر پیش‌بینی‌شده توسط مدل و مقادیر حقیقی را برای به‌روزرسانی پارامترها کمّی می‌سازد.

\item[گرادیان کاهشی تصادفی:] روشی بهینه‌سازی تکرارشونده که در هر گام، شیب تابع هدف را با استفاده از دسته‌ای کوچک از داده‌ها تخمین زده و پارامترها را تنظیم می‌نماید.
\end{descriptiveglossary}

% -----------------------------------------------------------------------------
% ۲. واژه‌نامه فارسی به انگلیسی
% -----------------------------------------------------------------------------
\clearpage
\thispagestyle{plain}
\setheader{واژه‌نامه فارسی به انگلیسی}
\chapter*{واژه‌نامه فارسی به انگلیسی}
\addcontentsline{toc}{chapter}{واژه‌نامه فارسی به انگلیسی}

\persianToEnglish{استخراج ویژگی}{Feature Extraction}
\persianToEnglish{الگوریتم بهینه‌سازی}{Optimization Algorithm}
\persianToEnglish{بازتولیدپذیری}{Reproducibility}
\persianToEnglish{برازش بیش از حد}{Overfitting}
\persianToEnglish{پردازش زبان طبیعی}{Natural Language Processing (NLP)}
\persianToEnglish{پرتو جستجو}{Beam Search}
\persianToEnglish{تابع فعال‌ساز}{Activation Function}
\persianToEnglish{تایپوگرافی دوزبانه}{Bidirectional Typography}
\persianToEnglish{تنظیم‌سازی}{Regularization}
\persianToEnglish{حروف‌چینی رایانه‌ای}{Computer Typesetting}
\persianToEnglish{خطای مربع میانگین}{Mean Squared Error (MSE)}
\persianToEnglish{خودرمزگذار}{Autoencoder}
\persianToEnglish{دسته‌بندی چندکلاسه}{Multi-class Classification}
\persianToEnglish{رابطه ریاضی}{Mathematical Equation}
\persianToEnglish{زی‌پرشین}{XePersian}
\persianToEnglish{سامانه بازشناسی تصویر}{Image Recognition System}
\persianToEnglish{سرصفحه}{Header}
\persianToEnglish{شبکه عصبی پیچشی}{Convolutional Neural Network (CNN)}
\persianToEnglish{شبکه عصبی عمیق}{Deep Neural Network (DNN)}
\persianToEnglish{فراسنج}{Hyperparameter}
\persianToEnglish{گرادیان کاهشی}{Gradient Descent}
\persianToEnglish{ماتریس درهم‌ریختگی}{Confusion Matrix}
\persianToEnglish{متن دوسویه}{Bidirectional Text}
\persianToEnglish{مدل زبان بزرگ}{Large Language Model (LLM)}
\persianToEnglish{نرخ یادگیری}{Learning Rate}
\persianToEnglish{هوش مصنوعی}{Artificial Intelligence (AI)}
\persianToEnglish{یادگیری تحت نظارت}{Supervised Learning}
\persianToEnglish{یادگیری تقویتی}{Reinforcement Learning}
\persianToEnglish{یادگیری عمیق}{Deep Learning}
\persianToEnglish{یادگیری ماشین}{Machine Learning}

% -----------------------------------------------------------------------------
% ۳. واژه‌نامه انگلیسی به فارسی
% -----------------------------------------------------------------------------
\clearpage
\thispagestyle{plain}
\setheader{واژه‌نامه انگلیسی به فارسی}
\chapter*{واژه‌نامه انگلیسی به فارسی}
\addcontentsline{toc}{chapter}{واژه‌نامه انگلیسی به فارسی}

\englishToPersian{Activation Function}{تابع فعال‌ساز}
\englishToPersian{Artificial Intelligence (AI)}{هوش مصنوعی}
\englishToPersian{Attention Mechanism}{سازوکار توجه}
\englishToPersian{Autoencoder}{خودرمزگذار}
\englishToPersian{Beam Search}{جستجوی پرتویی}
\englishToPersian{Bidirectional Text}{متن دوسویه}
\englishToPersian{Computer Typesetting}{حروف‌چینی رایانه‌ای}
\englishToPersian{Confusion Matrix}{ماتریس درهم‌ریختگی}
\englishToPersian{Convolutional Neural Network (CNN)}{شبکه عصبی پیچشی}
\englishToPersian{Deep Learning}{یادگیری عمیق}
\englishToPersian{Feature Extraction}{استخراج ویژگی}
\englishToPersian{Gradient Descent}{گرادیان کاهشی}
\englishToPersian{Header}{سرصفحه}
\englishToPersian{Hyperparameter}{فراسنج}
\englishToPersian{Image Recognition}{بازشناسی تصویر}
\englishToPersian{Large Language Model (LLM)}{مدل زبان بزرگ}
\englishToPersian{Learning Rate}{نرخ یادگیری}
\englishToPersian{Loss Function}{تابع زیان}
\englishToPersian{Machine Learning}{یادگیری ماشین}
\englishToPersian{Mean Squared Error (MSE)}{خطای مربع میانگین}
\englishToPersian{Natural Language Processing (NLP)}{پردازش زبان طبیعی}
\englishToPersian{Optimization Algorithm}{الگوریتم بهینه‌سازی}
\englishToPersian{Overfitting}{برازش بیش از حد}
\englishToPersian{Regularization}{تنظیم‌سازی}
\englishToPersian{Reinforcement Learning}{یادگیری تقویتی}
\englishToPersian{Reproducibility}{بازتولیدپذیری}
\englishToPersian{Supervised Learning}{یادگیری تحت نظارت}
\englishToPersian{Transformer}{ترنسفورمر}
\englishToPersian{XeLaTeX}{زی‌لاتک}
\englishToPersian{XePersian}{زی‌پرشین}
